\documentclass[letterpaper]{article}

\usepackage[letterpaper,margin=1in]{geometry}

% Compile from the repository root so that this path resolves to
% the pinned WSP submodule.
\input{wsp/testing/test-case-library.tex}

\begin{document}

\def\TCID{%
  TC-0005
}

\def\TCTitle{%
  Toolbar button rules
}

\def\TCPurpose{%
  Verify the rules of toolbar buttons: valid styles and values, what a press
  does for each style, the drop-down rule and the saved layout text.
}

\def\TCPriority{%
  \TCPriorityHigh
}

\def\TCRequirementRef{%
  REQ-0012, REQ-0020, REQ-0022
}

\def\TCDesignRef{%
  Not applicable; the rules are in the named header.
}

\def\TCFeatureRef{%
  \texttt{include/wfc/mscomctl/toolbar.hpp} and
  \texttt{tests/toolbar\_test.cpp} and the runner
  \texttt{tests/tc-0005-toolbar.ps1}
}

\def\TCTestTechnique{%
  Equivalence partitioning and boundary value analysis of the rule functions
  (ISO/IEC/IEEE 29119-4 specification-based techniques).
}

\def\TCPreconditions{%
  The library is built standalone (\texttt{cmake -S . -B out/build}, then
  \texttt{cmake --build out/build --config Debug}) with
  \texttt{WFC\_BUILD\_TESTS} on; CTest runs the test as
  \texttt{TC-0005-toolbar}.
}

\def\TCEnvironment{%
  Any target the library builds for (Windows x86, x64 and ARM64 with MSVC,
  Linux through Winelib); local runs use the x64 Debug configuration of the
  standalone build. The test needs no window, registry or network.
}

\def\TCAssumptions{%
  The rule header under test is the one built into the library; the test is
  compiled from the same source revision.
}

\def\TCInputData{%
  Styles -1 to 6, values -1 to 2, the six button styles, and the layouts
  \texttt{new=1;open=0;}, \texttt{a=1;junk;b=0} and the empty text.
}

\def\TCInitialState{%
  A fresh process; the test keeps no state between runs and holds no files or
  handles.
}

\def\TCProcedure{%
  \begin{enumerate}
    \item Run \texttt{mscomctl\_toolbar\_test} (CTest
      \texttt{TC-0005-toolbar}).
    \item Ask whether each style and value is valid, and what a press does for
      each style, enabled and disabled.
    \item Save a two-button layout to text, read it back, and read a text with
      a piece that has no \texttt{=} and no final \texttt{;}, and an empty
      text.
  \end{enumerate}
}

\def\TCExpectedResult{%
  \begin{enumerate}
    \item Styles 0 to 5 and values 0 and 1 are valid, others are not.
    \item A disabled button, a separator and a placeholder ignore a press;
      default and drop-down buttons only click; a check button toggles; a
      group button presses its group; \texttt{Value = tbrPressed} presses only
      a group button; only an enabled drop-down button drops down.
    \item The saved text is \texttt{new=1;open=0;} and reads back to the same
      two keys and visibility flags; a piece without \texttt{=} is skipped and
      the last piece needs no \texttt{;}; the empty text has no entries.
    \item The program prints \texttt{toolbar: ok} and exits with code 0.
  \end{enumerate}
}

\def\TCPostConditions{%
  The process ends; nothing is written outside CTest's own log.
}

\TCTemplate

\end{document}
